\section{Discussion}
\label{sec:discussion}

\textbf{Sources of Performance Gain.} The controlled ablation study reframes the
training design space presented in Section~\ref{sec:objectives}. Focal Loss and
many-hot supervision do not produce consistent practical gains over the BCE control, whereas
post-processing changes the operating point: Gaussian smoothing changes
boundary decisions, and temperature scaling improves probability quality
without changing the ranking. For a frozen-feature pipeline, this ordering
is practically important. The post-processing search runs on cached logits
and can be repeated for a deployment domain without retraining the
representation.

\textbf{Limitations of Aggregate Calibration Metrics.} After temperature scaling, the
equal-width ECE falls to 0.002214, yet the class-balanced ECE remains
0.070405. The gap follows from boundary sparsity: non-boundary frames
dominate aggregate confidence bins, so a low aggregate score can coexist
with poorer calibration on the rare positive class. Reporting only
aggregate ECE would therefore overstate calibration quality for SBD and
likely for other sparse temporal-event tasks. Event-conditional calibration,
such as fitting a separate temperature on boundary-adjacent frames, represents a
promising direction for future research, but it requires an independent calibration split that the
present protocol does not reserve.

\textbf{Analysis of Uncertainty Sources.} The two uncertainty analyses answer
complementary questions. Paired video bootstrap intervals
(Table~\ref{tab:paired_delta}) estimate how much the measured deltas depend
on which videos are sampled. The three-seed study estimates the sensitivity
of the controlled pipeline to head-training randomness. In our runs, the
seed-induced standard deviations are below 0.5 percentage points in F1
across all three datasets and are much smaller than the paired bootstrap
interval widths. The remaining unquantified term is protocol variability,
because
the frozen five-fold selection is itself a function of the validation
videos.

\textbf{Precision Bottleneck on ClipShots.} The transition-type
breakdown shows gradual-transition recall (0.8089) trailing cut recall
(0.8991), but the dominant loss is aggregate precision
(\PaperDeployClipPrecision): open-world content produces false-positive
peaks that survive smoothing. A frozen representation cannot adapt to this
distribution shift by construction, so threshold or temperature changes can
only trade recall for precision along a fixed curve. Closing the gap likely
requires domain-adaptive components, either lightweight adapters on the
frozen backbone or domain-specific post-processing selection, rather than
further tuning of the present loss-design space.
